Integrated terrestrial and low Earth orbit (LEO) satellite networks give an airborne or remote terminal several paths with very different behavior. A 5G link has low and fairly stable latency but limited coverage. Starlink and OneWeb links have global coverage but their round-trip time (RTT), jitter, and loss change quickly with satellite handovers, scheduler reconfigurations, and load~\cite{michel2022first,mohan2024multifaceted}. The Year 2 work of this project built a Multipath QUIC (MPQUIC) \emph{liquid tunnel} that combines these paths and showed that a scheduler fed with short-horizon forecasts of path state loses fewer video frames than a reactive minRTT scheduler~\cite{deconinck2017multipath}. The quality of those forecasts is therefore a direct input to the quality of the scheduling decisions.

The Year 2 Modeling Dynamics report compared eight forecasting models on six links (Starlink, OneWeb, and 5G, each in downlink and uplink) and selected a GRU sequence-to-sequence model with additive attention (Seq2Seq-GRU) as the primary model~\cite{cho2014learning,bahdanau2015neural}. This Year 3 report starts from that model and asks three questions.
\begin{enumerate}[leftmargin=*]
  \item \textbf{Are the Year 2 numbers right?} We re-implemented the evaluation, reproduced the Year 2 results, and audited the data and the metrics. Section~\ref{sec:audit} lists six issues. One of them, an interpolation step in the Starlink parser, leaks part of the next RTT sample into the model input. Another one makes the Year 2 interval coverage equal to 95\% by construction.
  \item \textbf{Can the forecasts be made better?} Section~\ref{sec:methods} describes the Year 3 changes: a causal data pipeline, a residual decoder anchored on the last observation, engineered context features, and, most importantly, distributional output heads that predict a mean and five quantiles per target and horizon.
  \item \textbf{Can the uncertainty be trusted?} The scheduler needs to know how wrong a forecast may be. We calibrate intervals with split conformal prediction on held-out sessions~\cite{lei2018distribution,romano2019conformalized} and report the coverage that is actually achieved on unseen test sessions.
\end{enumerate}

\paragraph{Related work.} Measurement studies explain why LEO link state changes quickly. Starlink reassigns users to satellites in globally synchronized 15\,s slots, which produces latency jumps~\cite{mohan2024multifaceted,tanveer2023making}. Its throughput varies from sub-second to diurnal scales~\cite{garcia2023multitimescale}. OneWeb latency changes at each satellite and gateway handover~\cite{zhao2025measuring}, which an earlier study on our testbed also observed~\cite{perrin2025measurement}. Learned prediction of link performance is established for cellular networks~\cite{narayanan2020lumos5g,mei2019realtime}. At sub-second resolution, however, simple predictors are hard to beat~\cite{basit2024predictability}. For LEO networks, Starlink throughput and latency have been predicted at coarse resolution~\cite{liu2025vivisecting,benghe2026horizon}. A concurrent preprint forecasts LEO link state with quantile bands~\cite{chi2026reading}. On the transport side, MPQUIC schedulers act on smoothed RTT estimates~\cite{paasch2014experimental,lim2017ecf}. Measured multipath sessions over Starlink and cellular paths suffer from the resulting mismatch~\cite{baltaci2023multipath}. Handover-aware or forecast-driven schedulers have been proposed~\cite{yang2024mams,maglione2025intelligent}. None of these works checks the causality of the input alignment or calibrates forecast uncertainty on unseen sessions, which are the two gaps this report addresses.

All methods are compared on the same session-level split as Year 2, with three training seeds per neural model, against four non-neural references (persistence, a one-second seasonal repeat, and two ridge regressions). Every number in this report is produced by the scripts listed in Appendix~\ref{app:repro} and is stored in \texttt{modeling/\allowbreak{}results\_\allowbreak{}y3/\allowbreak{}}.
